% ADDED figure (2026-11-14): Theorem §C8.12.5, the four-point (seagull) vertex of scalar QED, 2ie^2 g^{mu nu} (the user's PHY 513 notes, Appendix B, Derivation "Scalar QED"; Yu (11.495)). Generated by qft-C/tools/rulesL/figs/make_figs.sh
\tikzset{
  fdl/.style={line width=0.8pt, black},
  fdx/.style={circle, fill=black, inner sep=0pt, minimum size=3.6pt},
  fdv/.style={circle, fill=figB, draw=figB, inner sep=0pt, minimum size=5pt},
  fdlab/.style={font=\scriptsize, inner sep=2.6pt},
  fdmom/.style={->, >=Stealth, figO, line width=0.6pt},
}
\tikzset{
  fda/.style={fdl, postaction={decorate}, decoration={markings, mark=at position 0.55 with {\arrow{Stealth[length=5pt]}}}},
  fdS/.style={font=\small},
}
\tikzset{
  fline/.style={line width=0.8pt, black, postaction={decorate}, decoration={markings, mark=at position 0.56 with {\arrow{Stealth[length=5.5pt, width=4.5pt]}}}},
  sline/.style={line width=0.8pt, black, dash pattern=on 3.5pt off 2.2pt},
  blob/.style={draw, line width=0.8pt, fill=figB!12},
}
\tikzset{
  pline/.style={line width=0.8pt, black, decorate, decoration={snake, amplitude=1.6pt, segment length=5.2pt}},
}
\begin{tikzpicture}[font=\small]
  \tikzset{cx/.style={line width=0.8pt, black, dash pattern=on 3.5pt off 2.2pt, postaction={decorate}, decoration={markings, mark=at position 0.58 with {\arrow{Stealth[length=5.5pt, width=4.5pt]}}}},
           asg/.style={line width=0.8pt, figO, -{Stealth[length=4.5pt]}}}
  \begin{scope}
    \draw[cx] (-1.3,-1.2)--(0,0); \draw[cx] (0,0)--(-1.3,1.2);
    \draw[pline] (0,0)--(1.4,1.2); \draw[pline] (0,0)--(1.4,-1.2); \node[fdv] at (0,0){};
    \node[fdlab,left] at (-1.3,-1.2){$\phi$}; \node[fdlab,left] at (-1.3,1.2){$\phi^\ast$};
    \node[fdlab,right] at (1.4,1.2){$A_\mu$}; \node[fdlab,right] at (1.4,-1.2){$A_\nu$};
    \node[fdS, anchor=west] at (2.3,0) {$= 2ie^2\,g^{\mu\nu}$};
    \node[fdS] at (0.6,-2.0) {(a) the vertex};
  \end{scope}
  \begin{scope}[xshift=7.6cm]
    \node[fdS] at (1.9,-2.0) {(b) the factor $2$: two photon lines on two $A$'s};
    \foreach \x/\a/\b/\t in {0/1/2/{$g^{\mu\nu}$}, 2.6/2/1/{$g^{\nu\mu}$}}{
      \begin{scope}[xshift=\x cm]
        \node[draw, circle, line width=0.7pt, inner sep=0pt, minimum size=5.5mm] (m) at (0,1.0) {$\mu$}; \node[draw, circle, line width=0.7pt, inner sep=0pt, minimum size=5.5mm] (n) at (1.2,1.0) {$\nu$};
        \node[draw, line width=0.7pt, rounded corners=2pt, fill=figB!10, minimum width=7.5mm, minimum height=5.5mm] (g1) at (0,-0.4) {$A_\rho$}; \node[draw, line width=0.7pt, rounded corners=2pt, fill=figB!10, minimum width=7.5mm, minimum height=5.5mm] (g2) at (1.2,-0.4) {$A_\sigma$};
        \draw[asg] (m) -- (g\a); \draw[asg] (n) -- (g\b);
        \node[fdlab] at (0.6,-1.1) {\t};
      \end{scope}}
    \node[fdlab] at (1.9,1.95) {the term $e^2g^{\rho\sigma}A_\rho A_\sigma\,\phi^\ast\phi$ and the lines $\mu$, $\nu$};
  \end{scope}
\end{tikzpicture}
